% Options for packages loaded elsewhere
\PassOptionsToPackage{unicode}{hyperref}
\PassOptionsToPackage{hyphens}{url}
\PassOptionsToPackage{dvipsnames,svgnames,x11names}{xcolor}
\documentclass[
  10pt,
  fontset=fandol]{ctexart}
\usepackage{xcolor}
\usepackage[margin=16mm]{geometry}
\usepackage{amsmath,amssymb}
\setcounter{secnumdepth}{-\maxdimen} % remove section numbering
\usepackage{iftex}
\ifPDFTeX
  \usepackage[T1]{fontenc}
  \usepackage[utf8]{inputenc}
  \usepackage{textcomp} % provide euro and other symbols
\else % if luatex or xetex
  \usepackage{unicode-math} % this also loads fontspec
  \defaultfontfeatures{Scale=MatchLowercase}
  \defaultfontfeatures[\rmfamily]{Ligatures=TeX,Scale=1}
\fi
\usepackage{lmodern}
\ifPDFTeX\else
  % xetex/luatex font selection
\fi
% Use upquote if available, for straight quotes in verbatim environments
\IfFileExists{upquote.sty}{\usepackage{upquote}}{}
\IfFileExists{microtype.sty}{% use microtype if available
  \usepackage[]{microtype}
  \UseMicrotypeSet[protrusion]{basicmath} % disable protrusion for tt fonts
}{}
\makeatletter
\@ifundefined{KOMAClassName}{% if non-KOMA class
  \IfFileExists{parskip.sty}{%
    \usepackage{parskip}
  }{% else
    \setlength{\parindent}{0pt}
    \setlength{\parskip}{6pt plus 2pt minus 1pt}}
}{% if KOMA class
  \KOMAoptions{parskip=half}}
\makeatother
\usepackage{longtable,booktabs,array}
\newcounter{none} % for unnumbered tables
\usepackage{calc} % for calculating minipage widths
% Correct order of tables after \paragraph or \subparagraph
\usepackage{etoolbox}
\makeatletter
\patchcmd\longtable{\par}{\if@noskipsec\mbox{}\fi\par}{}{}
\makeatother
% Allow footnotes in longtable head/foot
\IfFileExists{footnotehyper.sty}{\usepackage{footnotehyper}}{\usepackage{footnote}}
\makesavenoteenv{longtable}
\setlength{\emergencystretch}{3em} % prevent overfull lines
\providecommand{\tightlist}{%
  \setlength{\itemsep}{0pt}\setlength{\parskip}{0pt}}
\usepackage{amsmath,amssymb,mathtools,needspace}
\xeCJKDeclareCharClass{CJK}{"2032,"0400->"04FF,"2160->"216B,"2460->"2473,"25A0,"25C7,"25CB,"25CF}
\IfFontExistsTF{Arial Unicode MS}{\xeCJKsetup{AutoFallBack=true}\setCJKfallbackfamilyfont{\CJKrmdefault}{Arial Unicode MS}}{}
\setcounter{secnumdepth}{0}
\setlength{\emergencystretch}{2em}
\allowdisplaybreaks
\usepackage{bookmark}
\IfFileExists{xurl.sty}{\usepackage{xurl}}{} % add URL line breaks if available
\urlstyle{same}
\hypersetup{
  pdftitle={Adams 显式格式},
  colorlinks=true,
  linkcolor={Maroon},
  filecolor={Maroon},
  citecolor={Blue},
  urlcolor={Blue},
  pdfcreator={LaTeX via pandoc}}

\title{Adams 显式格式}
\author{}
\date{}

\begin{document}
\maketitle

\subsubsection{5.5.2 Adams 显式格式}\label{adams-ux663eux5f0fux683cux5f0f}

运用插值方法的关键，在于选取合适的插值节点。记 \(f_k=f(x_k,y_k)\)，先用 \(r+1\) 个数据点 \((x_n,\allowbreak f_n),\allowbreak (x_{n-1},\allowbreak f_{n-1}),\allowbreak \cdots,\allowbreak (x_{n-r},\allowbreak f_{n-r})\) 构造插值多项式 \(P_r(x)\)，注意到这里插值节点 \(x_n,x_{n-1},\cdots,x_{n-r}\) 等距，运用 Newton 后插公式（见 2.4.2 节）可写出

\[
P_r(x_n+th)=\sum_{j=0}^r(-1)^j\binom{-t}{j}\Delta^j f_{n-j},
\]

其中 \(t=\dfrac{x-x_n}{h}\)，\(\Delta^j\) 表示 \(j\) 阶向前差分，而

\[
\binom sj=\frac{s(s-1)\cdots(s-j+1)}{j!}.
\]

将 \(P_r(x)\) 的上述表达式代入式（5.5.2），即得下列 \textbf{Adams 显式格式}

\[
y_{n+1}=y_n+h\sum_{j=0}^r\alpha_j\Delta^j f_{n-j},
\tag{5.5.3}
\]

其中 \(\alpha_j\) 为不依赖于 \(n\) 和 \(r\) 的系数，\(\alpha_j=(-1)^j\int_0^1\binom{-t}{j}\,\mathrm dt\)。它的前几个值如表 5.6 所示。

\textbf{表 5.6}

{\def\LTcaptype{none} % do not increment counter
\begin{longtable}[]{@{}lrrrr@{}}
\toprule\noalign{}
\(j\) & 0 & 1 & 2 & 3 \\
\midrule\noalign{}
\endhead
\bottomrule\noalign{}
\endlastfoot
\(\alpha_j\) & 1 & \(1/2\) & \(5/12\) & \(3/8\) \\
\end{longtable}
}

实际计算时，将格式（5.5.3）中的差分展开往往是方便的，利用差分展开式 \(\Delta^j f_{n-j}=\sum_{i=0}^r(-1)^i\binom ji f_{n-j}\)，可将它改写为

\[
y_{n+1}=y_n+h\sum_{i=0}^r\beta_{ri}f_{n-i},\qquad
\beta_{ri}=(-1)^i\sum_{j=i}^r\binom ji\alpha_j.
\tag{5.5.4}
\]

\begin{quote}
校注（agent 补充）：原图中上述未编号``差分展开式''的求和上限印作 \(r\)，求和项末尾印作 \(f_{n-j}\)，这里均保留原文。末尾下标疑为原书排印错误，通常应为 \(f_{n-i}\)；例如 \(j=r=1\) 时，原印刷式右边为零，而左边为 \(f_n-f_{n-1}\)。按通常写法求和上限为 \(j\)；若约定 \(\binom ji=0\)（\(i>j\)），则上限写 \(r\) 也可。此疑误已对照原页局部图，并非 OCR 替换。
\end{quote}

\href{../../source-images/pdf-139.jpeg}{原页}

这里，系数 \(\beta_{ri}\) 与 \(r\) 的定值有关，其具体数值如表 5.7 所示。式（5.5.4）是含有参数 \(r\) 的一族格式，\(r+1\) 为格式的步数。特别地，一步显式 Adams 格式（\(r=0\)）即为 Euler 格式。两步显式 Adams 格式（\(r=1\)）为

\[
y_{n+1}=y_n+\frac h2(3f_n-f_{n-1}),
\tag{5.5.5}
\]

\textbf{表 5.7}

{\def\LTcaptype{none} % do not increment counter
\begin{longtable}[]{@{}lrrrr@{}}
\toprule\noalign{}
\(i\) & 0 & 1 & 2 & 3 \\
\midrule\noalign{}
\endhead
\bottomrule\noalign{}
\endlastfoot
\(\beta_{0i}\) & 1 & & & \\
\(\beta_{1i}\) & \(3/2\) & \(-1/2\) & & \\
\(\beta_{2i}\) & \(23/12\) & \(-16/12\) & \(5/12\) & \\
\(\beta_{3i}\) & \(55/24\) & \(-59/24\) & \(37/24\) & \(-9/24\) \\
\end{longtable}
}

而四步显式 Adams 格式（\(r=3\)）则为

\[
y_{n+1}=y_n+\frac h{24}(55f_n-59f_{n-1}+37f_{n-2}-9f_{n-3}).
\tag{5.5.6}
\]

\subsection{本节覆盖原页的校注记录}\label{ux672cux8282ux8986ux76d6ux539fux9875ux7684ux6821ux6ce8ux8bb0ux5f55}

下列记录按原页关联，含同页相邻小节的校注；计算时同时读取上文校注和完整原页。

\begin{itemize}
\tightlist
\item
  \texttt{ch05b-E001}：原书 PDF 页 138
\end{itemize}

\end{document}
